\begin{abstract}
Deploying an edited representation requires more than showing that one probe
fails on one validation split: the released representation may face an
arbitrary downstream attacker and a shifted deployment population. We present
MOSAIC, Minimax-Optimized Source-Agnostic Invariant Channels, a software-only
method that jointly selects a stochastic representation-release channel and a
task decoder while certifying privacy and utility under a structured
pre-release shift. A single multinomial confidence event over fine-token laws
simultaneously covers the continuum of channels selected on that same table,
so the certificate incurs no candidate-count correction. We derive an exact
multiclass attacker envelope, a shift-transfer theorem combining common drift
with bounded differential contamination, a matching utility certificate, and
a no-free-lunch result for unrestricted differential shift. The resulting
finite problem is solved globally by decoder enumeration and branch-aware
mixed-integer linear optimization. The empirical claims in this draft are
populated only from a hash-locked confirmatory study and independently replayed
receipts.
\end{abstract}

\section{Introduction}

Representation editing is often evaluated by fitting a probe to a frozen
embedding and checking whether a designated source attribute remains
predictable. That evaluation leaves two deployment gaps. First, failure of a
particular probe does not prevent a different downstream rule from recovering
the attribute. Second, an edit chosen on a validation table can exploit
sampling noise, while the deployed population can change after certification.
These gaps are especially consequential when the released representation is a
public interface: the designer controls the release mechanism but not the
future attacker.

This paper studies a finite but expressive release model. A tokenizer fixed
without the certification fold maps an internal representation to a fine token
$C\in[K]$. A source-blind row-stochastic channel $M\in\mathrm{Stoch}(K,L)$
releases the only public token $Z$. For each task label, privacy is measured by
the balanced accuracy of the Bayes-optimal source attacker on $Z$, not by a
chosen probe family. The deployment shift occurs before the software-enforced
channel and decomposes into a common Markov transformation plus a bounded
fraction of source-specific contamination. This model permits large
source-blind drift while explicitly charging the component that can create new
source differences.

MOSAIC, or Minimax-Optimized Source-Agnostic Invariant Channels, turns this
model into a certifying release rule. Its first ingredient is an exact robust
attacker envelope built from simultaneous multinomial confidence regions over
the fine-token table. Because the confidence event is stated before applying
$M$, it covers every stochastic channel at once, including one optimized on the
same data. Its second ingredient is an exact channel capacity for arbitrary
differential shift together with a computable defect for common drift. The
same event also certifies a jointly selected task decoder. Finally, the
piecewise-linear certificate is optimized globally for finite output and
decoder alphabets, and the method abstains when no release satisfies the
registered privacy and utility contract.

The contribution is not the individual use of randomized channels,
distributional confidence regions, contraction coefficients, or
accept-or-abstain decisions, all of which have substantial prior literatures.
The technical claim is their specific combination: continuum-wide same-table
selection, an exact arbitrary-attacker metric, structured pre-release shift,
simultaneous task utility, and a globally checked optimizer. We also prove that
unrestricted source-specific shift forces a constant-row release and quantify
the resulting privacy--utility obstruction. A hash-locked experiment compares
MOSAIC with plug-in selection, held-out fixed-channel certification, finite
Learn-Then-Test certification, deterministic release, a shift-unaware
ablation, and an exact population oracle.

\section{Related Work}

Universal fair-representation certificates protect against downstream
classifiers by constraining statistical distances or Bayes distinguishability
on a finite representation. Randomized preprocessing and privacy-funnel work
optimize stochastic mappings under population privacy and utility constraints.
MOSAIC adopts those established objects but asks a different statistical
question: can the stochastic post-channel itself be selected from an
uncountable family on the table used to certify it? The answer comes from
placing the simultaneous confidence event on the sufficient fine-token laws
rather than testing channels one by one.

Risk-control and Learn-Then-Test procedures certify members of registered
candidate families and abstain when no candidate passes. They are essential
baselines and the closest decision-layer precedent. MOSAIC does not claim a
new multiple-testing principle. Its envelope instead covers all stochastic
post-channels pointwise on one event, avoiding a candidate-family correction
and permitting a global continuum optimizer. Work on Dobrushin contraction and
domain generalization likewise establishes that stochastic representations
can contract distribution shift. Here the contraction quantity enters an
exact source-inference capacity and is combined with a separate invariance
defect for common pre-release drift.

Robust fair learning and impossibility results show that arbitrary population
shift can destroy nontrivial fairness guarantees. Our shift assumption is
therefore explicit rather than implicit. The common transformation and retained
mass are shared across sources; only a bounded contamination fraction may vary
by source. We certify performance conditionally on membership in this class
and do not infer real-world membership from an ordinary benchmark split.

\section{Problem Setup}

Let $S\in[G]$ denote source, $Y\in[J]$ task label, and $C\in[K]$ the fine token.
For a fixed label, write $p_s(c)=\Pr(C=c\mid S=s,Y=y)$. The source-blind release
channel $M$ produces $Z\in[L]$. The balanced accuracy of the strongest
downstream source attacker is
\begin{equation}
 A(P,M)=\frac{1}{G}\sum_{z=1}^{L}\max_{s\in[G]}(p_sM)(z),
 \label{eq:attacker}
\end{equation}
with normalized advantage
$\operatorname{Adv}_G(A)=(GA-1)/(G-1)$. This is chance-normalized and ranges
from zero to one.

For each supported source--label stratum, the certification table supplies an
empirical law $\widehat p_{s,y}$ and radius $\epsilon_{s,y}$. The simultaneous
event
\begin{equation}
 \mathcal E=\bigcap_{s,y}
 \{\|\widehat p_{s,y}-p_{s,y}\|_1\leq\epsilon_{s,y}\}
 \label{eq:event}
\end{equation}
has probability at least $1-\delta$ after a fixed allocation across strata. We
use the Weissman multinomial radius in the confirmation, although any valid
simultaneous region can replace it. The tokenizer and fine alphabet must be
fixed without these observations.

The external law for label $y$ belongs to the structured class
\begin{equation}
 q_s=t_y p_sT_y+(1-t_y)r_s,
 \qquad t_y\geq 1-\eta_y,
 \label{eq:shift}
\end{equation}
where $T_y$ is a source-independent fine-token channel from a registered
convex uncertainty set and each residual $r_s$ is arbitrary. Both components
pass through the same release channel $M$. This placement is crucial: a
constant-row software channel cannot leak information that appears before it.

\section{Adaptive Universal-Attacker Certificate}

For a bounded vector $v$, define the exact confidence-set support function
\begin{equation}
 \Phi(\widehat p,\epsilon;v)=
 \max_{p\in\Delta_K:\|p-\widehat p\|_1\leq\epsilon}p^\top v.
 \label{eq:phi}
\end{equation}
It is obtained by moving at most $\epsilon/2$ mass from low to high
coordinates of $v$, or from an equivalent linear-program dual. For attacker
assignment $a\in[G]^L$, let
$v_{M,a,s}(c)=\sum_zM(c,z)\mathbf 1\{a(z)=s\}$. We then define
\begin{equation}
 \overline A(\widehat P,M)=\frac{1}{G}\max_{a\in[G]^L}
 \sum_s\Phi(\widehat p_s,\epsilon_s;v_{M,a,s}).
 \label{eq:abar}
\end{equation}

\begin{theorem}[Adaptive exact envelope]
On $\mathcal E$, simultaneously for every row-stochastic $M$, including a
channel selected as an arbitrary function of the same table,
$A(P,M)\leq\overline A(\widehat P,M)$. For fixed $M$, the right-hand side is the
exact supremum over the product of the stated $\ell_1$ confidence sets.
\end{theorem}

The proof writes the Bayes attacker as a maximum over deterministic assignments.
For a fixed assignment the source confidence sets separate, giving one support
function per source; all maxima are finite and commute. The implication is
pointwise over the complete stochastic-channel space, which is why adaptive
selection needs no channel-count correction.

\section{Certification under Pre-Release Shift}

Define the arbitrary differential-shift capacity
\begin{equation}
 \operatorname{Cap}_G(M)=\frac{1}{G}\max_{a\in[G]^L}
 \sum_s\max_c v_{M,a,s}(c).
 \label{eq:capacity}
\end{equation}
This equals $\sup_{r_1,\ldots,r_G}A(R,M)$. For $G=2$, its normalized value is
the Dobrushin coefficient, the largest total variation distance between two
rows of $M$.

For common transformation $T$, define
$\rho_T(M)=\max_c\operatorname{TV}((TM)(c,\cdot),M(c,\cdot))$ and take the
maximum over registered extremes. The retained common component obeys
$A(P,TM)\leq\min\{\operatorname{Cap}_G(M),A(P,M)+\rho_T(M)\}$. Combining this
transfer with the contamination decomposition gives the main certificate.

\begin{theorem}[Structured-shift certificate]
On $\mathcal E$, every adaptively selected $M$ and every external law in
Eq.~\eqref{eq:shift} satisfy
\begin{equation}
 A(Q,M)\leq(1-\eta_y)B_y(M)+\eta_y\operatorname{Cap}_G(M),
 \label{eq:mainbound}
\end{equation}
where
$B_y(M)=\min\{\operatorname{Cap}_G(M),
\overline A_y(\widehat P,M)+\rho_{\mathcal T_y}(M)\}$.
\end{theorem}

Bayes accuracy is convex in the source experiment. The common component is
bounded by invariance transfer and the residual by channel capacity. Since the
common bound cannot exceed capacity, the worst retained mass occurs at
$t_y=1-\eta_y$. When $\eta_y=0$ and $TM=M$, the external certificate equals the
reference certificate despite common drift. When $\eta_y=1$, it reduces to the
exact distribution-free capacity rather than an assumed external validation
event.

For decoder $g:[L]\to[J]$, let
$u_y(c)=\sum_zM(c,z)\mathbf 1\{g(z)\neq y\}$. Replacing attacker scores by this
bounded loss yields a simultaneous conditional-error certificate. Its common
defect is the maximum positive increase in $u_y$ under $T$, and its differential
capacity is $\max_cu_y(c)$. Thus $M$, $g$, privacy, and worst-stratum utility can
all be selected on the same event.

\begin{theorem}[No free lunch]
$\operatorname{Cap}_G(M)=1/G$ if and only if all rows of $M$ are identical.
For two sources, if two fine tokens carry different task labels and incur
row-wise decoder errors at most $e_0$ and $e_1$, then the normalized capacity is
at least $1-e_0-e_1$.
\end{theorem}

Consequently, unrestricted source-specific shift permits exact erasure only by
making the public token independent of the fine token. A bounded $\eta_y$ is
therefore an identifiable assumption of the deployment contract, not a hidden
technical convenience.

\section{Global Optimization and Abstention}

For fixed output size, decoder, finite transform extremes, and attacker
assignments, every support-function dual and certificate branch has a linear
epigraph. Binary variables choose between capacity and coupled-shift branches;
the implementation enumerates all finite decoders and solves one mixed-integer
linear program per decoder. A release is accepted only when the solver reports
zero numerical MIP gap, its dual bound matches, all primal residuals pass, and
independently recomputed certificates agree with the objective. Otherwise the
rule returns an explicit abstention.

The optimizer is checked against exhaustive stochastic-channel grids on random
small problems and against exact population-risk enumeration. A constructed
witness also tests whether stochastic release is genuinely necessary: its
globally certified stochastic error is compared with every deterministic
channel. These checks validate the finite encoding but do not replace the
algebraic certificate.

The locked analysis also predicts abstention as sample size changes. A
finite-sample envelope follows by enlarging population-centered confidence
balls: whenever the enlarged problem remains feasible, a Weissman union bound
controls the probability that the empirical optimizer abstains. Because this
bound can be loose near a decision boundary, we separately preregister a local
delta-method power curve. Under a unique stable branch and strong regularity of
the parametric linear program, centered simplex finite differences estimate the
value gradient and the stratified multinomial central limit theorem predicts
the deployment probability. A two-step gradient-stability check is required;
the approximation is rejected if it fails.

\section{Hash-Locked Evaluation}

The synthetic confirmation fixes two source--label populations, two structured
shift regimes, eight deployment rules, five possible per-stratum sample sizes,
and 1,000 independent tables per registered cell. Pilot seeds are permanently
excluded. Before confirmatory outcomes, the complete protocol, method code,
theorem text, evaluators, pilot history, pass criteria, and theory predictions
are SHA-256 locked and pushed to a public branch. Every selected channel is
graded by a separately implemented exact population evaluator that enumerates
common-transform extremes, attacker assignments, and residual simplex
vertices. The confirmation stores every channel, decoder, seed, decision,
confidence-event marker, and exact risk, then an independent script replays all
rows and aggregate intervals.

The primary safety contrast compares the plug-in continuum rule with MOSAIC at
a near-boundary shift contract. The primary usefulness contrast compares safe
retention against a held-out fixed-channel certificate, a finite-grid
Learn-Then-Test rule with simultaneous one-sided bounds, and deterministic
MOSAIC. The final gate additionally requires the preregistered theory curve to
track deployment across sample sizes. All methods and cells remain reportable
if any gate fails.

\section{Locked Synthetic Results}

All preregistered gates pass. Figure~\ref{fig:confirmation} reports every
implementable rule in the primary cells; ``fallback'' is MOSAIC's generic
capacity certificate. Independent replay recomputes all 64,000 method receipts,
128,000 population risks, decisions, and aggregate intervals with zero
mismatch.

\begin{figure*}[t]
  \centering
  \includegraphics[width=\textwidth]{figures/figure2_mosaic_confirmation.pdf}
  \caption{Hash-locked confirmation with 1,000 independent certification
  tables per cell.  Error bars are exact 95\% Clopper--Pearson intervals.
  (a) At the safety boundary, plug-in continuum selection falsely accepts
  42.1\% of tables; MOSAIC abstains with no false acceptance. Ignoring shift is
  worse. (b) In the retention cell, the fallback safely deploys on 56.3\%,
  versus 2.5\% for held-out certification, 3.3\% for finite LTT, and 0\% for
  deterministic MOSAIC. (c) Deployment follows the preregistered theory curve
  with 0.20 percentage-point mean absolute error. Exact 95\%
  Clopper--Pearson intervals are shown.}
  \label{fig:confirmation}
\end{figure*}

The paired safety difference is 42.1 points (95\% bootstrap interval
$[39.0,45.2]$), with 421 plug-in-only violations and no reverse discordance.
Across 8,000 fallback certification tables, no accepted release violates the
exact population contract. In the retention cell, fallback safe deployment
exceeds held-out, finite-LTT, and deterministic rates by 53.8, 53.0, and 56.3
points. No deterministic channel is feasible. These finite-run facts support,
but do not replace, the $1-\delta$ guarantee.

A separately locked 5,000-table refinement evaluates exact shift geometry on
the same tables. At $n=125$, transform-exact certification safely deploys on
53.0\% versus 0.3\% for fallback; at $n=250$, the rates are 99.9\% and 57.8\%.
The exact objective is no larger on all 5,000 pairs, as theory requires, and all
10,000 method-cell decisions have zero false acceptances. Full cell tables,
threshold sensitivity, amendment history, and audits are in the supplement.


\section{Limitations and Scope}

The statistical guarantee is conditional on the shift class in
Eq.~\eqref{eq:shift}. An ordinary benchmark split cannot prove that a real
hospital, text domain, or sensor environment belongs to that class. Such a
claim requires an engineering specification or separate bridge data with its
own confidence allocation. The current finite-token theorem also assumes that
the tokenizer and alphabet are fixed independently of the certification table.
Large alphabets may require structured confidence regions or column generation
rather than full assignment enumeration.

MOSAIC certifies source inference from the released token, not every possible
harm associated with representation use. It does not protect information that
is released through a separate interface, and it cannot recover utility that
the privacy contract makes information-theoretically impossible. Finally, the
novelty claim concerns the complete adaptive and shift-aware construction; its
classical ingredients must remain attributed to their respective literatures.

\section{Conclusion}

MOSAIC treats representation editing as release-mechanism design rather than a
probe contest. By certifying fine-token laws before a stochastic post-channel,
it can optimize over a continuum on the same table while retaining a single
finite-sample event. The structured shift theorem separates common drift that
the channel can absorb from differential drift that its capacity must pay for,
and the no-free-lunch result states when abstention is unavoidable. The locked
evaluation determines whether those guarantees retain enough safe utility to
justify the method in practice.
